\documentclass[11pt]{article}

\usepackage[margin=0.86in]{geometry}
\usepackage{amsmath}
\usepackage{amssymb}
\usepackage{booktabs}
\usepackage{microtype}
\usepackage{array}
\usepackage{longtable}
\usepackage{url}
\usepackage{xcolor}
\usepackage[hidelinks]{hyperref}

\definecolor{passgreen}{HTML}{047857}
\definecolor{failred}{HTML}{B91C1C}
\definecolor{resultblue}{HTML}{1D4ED8}
\newcommand{\pass}{\textcolor{passgreen}{\textbf{PASS}}}
\newcommand{\fail}{\textcolor{failred}{\textbf{FAIL}}}
\newcommand{\code}[1]{\texttt{#1}}

\title{\textbf{Deterministic ``Neurons'' Inside an Existing Transformer MLP}\\
\large A Three-Seed Causal Study of Natural-Language Addition,
Routing, Preservation, and Registered State in Qwen2.5-0.5B}
\author{Josiah Wilson\\Independent Researcher}
\date{July 26, 2026}

\begin{document}
\maketitle

\begin{abstract}
We test whether a pretrained language model can learn to use a deterministic
calculator that occupies a small activation subspace inside an existing
transformer MLP, rather than invoking an external tool or replacing an answer
after generation. We modified frozen Qwen2.5-0.5B-Instruct at decoder layer 16.
Twenty-eight of its existing 4,864 MLP coordinates form a typed interface:
sixteen learned route/operand/digit coordinates feed a frozen decimal
ripple-carry adder, and twelve deterministic result coordinates return through
learned columns of the ordinary MLP down projection. The residual and MLP
widths do not grow. The interface has 25,088 learned weights
($0.0051\%$ of the 494-million-parameter base); the calculator has zero.

Development was explicitly adaptive and all failed audits were retained. An
initial shared three-seed holdout produced 58/60 exact additions per seed and
a 55/60 paired result-ablation loss per seed, but falsely routed 17/180
adversarial negatives. Retraining only the 1,792 route-row weights yielded
0/360 false routes and 360/360 token-exact preservation across two subsequent
exact-string-disjoint audits. A deterministic response-local operand register
then eliminated a sequence-time state-loss failure.

In the final frozen register audit, three independently learned interfaces
answered 173/180 additions exactly (96.11\%), including 90/90 word problems.
They recovered 174/180 exact operand registers; all 174 remained stable and
produced exact calculator digit-plus-end trajectories. Calculator-result
ablation left 9/180 correct, a paired causal loss on 165/180 examples. All
180/180 adversarial negatives remained route-off and token-identical to
untouched Qwen. The compound protocol nevertheless failed: mean exact accuracy
was 57.67/60 against a frozen 58/60 threshold, and one seed decoded exact
calculator symbols for $0+0$ as the token ``I.'' The result strongly supports
the narrow causal implant hypothesis while locating remaining failures in
learned input/output interfaces. Route, operand, and position state remain
generation-runtime microcircuit state; arbitrary recurrent multi-call
reasoning is not established.
\end{abstract}

\section*{Technical summary}

\begin{center}
\small
\begin{tabular}{p{0.48\linewidth}r}
\toprule
Final frozen audit outcome & Result\\
\midrule
Exact additions & 173/180 (96.11\%)\\
Exact first-step operand registers & 174/180 (96.67\%)\\
Exact calculator trajectories & 174/180 (96.67\%)\\
Exact natural-language word problems & 90/90\\
Exact after calculator-result ablation & 9/180 (5.00\%)\\
Paired correct-to-incorrect ablations & 165/180\\
Stable registers given exact first-step operands & 174/174\\
False routes on adversarial negatives & 0/180\\
Token-exact negative preservation & 180/180\\
\bottomrule
\end{tabular}
\end{center}

\noindent\textbf{Narrow mechanistic verdict.}
The in-place path works. Learned Qwen activations supplied typed operands to a
fixed calculator; exact internal symbols were causally necessary for most
answers; and ordinary frozen downstream layers usually decoded those symbols
correctly.

\medskip
\noindent\textbf{Compound protocol verdict.}
\fail. Three of five final engineering gates passed. Mean addition accuracy
missed its frozen threshold by one aggregate answer, and one exact internal
trajectory was misdecoded downstream.

\medskip
\noindent\textbf{Principal boundary.}
This is a supervised single-addition activation implant with runtime-managed
state. It is not evidence for unlimited calculator calls, spontaneous
discovery during pretraining, four operations, residual-native recurrent
registers, or arbitrary-prompt safety.

\section{Motivation and prior work}

Language models are flexible probabilistic sequence predictors, while
arithmetic algorithms are rigid state-transition systems. Learned arithmetic
architectures can improve extrapolation, but neural models still often fail
outside trained lengths or prompt distributions
\cite{trask2018,jelassi2023,cho2024}. Compiled transformers and neural
algorithmic reasoning instead ask how explicit algorithmic structure can
coexist with learned representations \cite{lindner2023,bounsi2024}.

The closest direct prior work is Dietz and Klakow's Integrated Gated
Calculator (IGC) \cite{dietz2025}. IGC freezes Llama 3.1 8B and learns
input, routing, and output mappings around a non-differentiable
four-operation calculator. It establishes prior art for a calculator
integrated into an LLM forward process. The present project does not claim the
broad integrated-calculator concept as novel.

The more specific hypothesis studied here was proposed by Josiah Wilson:
selected neuron-like coordinates in an otherwise ordinary network can be
assigned deterministic semantics, allowing surrounding learned computation to
discover when and how to use exact machinery. In the limiting version, a model
trained from scratch would experience the calculator coordinates like any
other neurons, except that their transfer function implements a known process.
This experiment asks a practical pretrained-model question:

\begin{quote}
Can a small frozen Qwen model learn to route ordinary language and distributed
number representations through a deterministic subspace inside one existing
MLP, then use the result through its ordinary downstream network while
preserving unrelated behavior?
\end{quote}

Earlier phases of this project tested output-adjacent residual bridges, an
internal block wrapper, natural-language routing with a fixed parser, a direct
comparison with an IGC-style learned parser, and a multi-stage neural
interface. Phase 7 is the first phase that replaces coordinates of an existing
Qwen MLP activation tensor itself.

\section{Architecture}

\subsection{Frozen pretrained model}

Every experiment uses \code{Qwen/Qwen2.5-0.5B-Instruct}, revision
\begin{sloppypar}
\code{7ae557604adf67be50417f59c2c2f167def9a775}
\end{sloppypar}
\cite{qwen2024}. The model has 494,032,768 pretrained parameters, 24 decoder
blocks, residual width 896, and MLP intermediate width 4,864. All pretrained
weights remain frozen. Generation is greedy.

\subsection{In-place activation ABI}

For hidden state $h$ at decoder layer 16 (zero-indexed), the original Qwen MLP
computes
\[
  y_{\mathrm{MLP}}
  = W_{\mathrm{down}}
    \left(\operatorname{SiLU}(W_{\mathrm{gate}}h)
    \odot W_{\mathrm{up}}h\right).
\]
A census selected 28 low-impact coordinates in the 4,864-wide intermediate
activation. Their original contribution is separated from the unselected
base computation. The implant assigns:

\begin{center}
\small
\begin{tabular}{lrr}
\toprule
Activation group & Coordinates & Learned weights\\
\midrule
Route: \code{OFF}, \code{ADD} & 2 & $2\times896=1,792$\\
Token role: none, operand A, operand B & 3 & $3\times896=2,688$\\
Digit type: 0--9, non-digit & 11 & $11\times896=9,856$\\
Result: 0--9, end, pad & 12 & $12\times896=10,752$\\
\midrule
Total & 28 & 25,088\\
\bottomrule
\end{tabular}
\end{center}

The first sixteen rows classify route, role, and digit from Qwen's native
layer-16 residuals. Role and digit predictions must agree, and digit confidence
must exceed 0.90, before a token enters an operand. Hard typed operands feed a
zero-parameter decimal ripple-carry cell. Its one-hot result symbol is mapped
through twelve learned down-projection columns into Qwen's residual stream.
The remaining seven decoder blocks, final normalization, and tied language
head are unchanged.

The calculator is not a post-generation correction. Its activation is injected
while the model computes the next token, and its result must survive the normal
downstream network.

\begin{center}
\fbox{\begin{minipage}{0.92\linewidth}
\centering
Frozen Qwen blocks 0--15\\[2pt]
$\downarrow$\\[2pt]
\textbf{Layer-16 MLP:}\\
ordinary 4,836 channels \quad $\parallel$ \quad
\begin{tabular}{c}
16 learned route/role/digit coordinates\\
$\downarrow$\\
frozen exact adder + registered state\\
$\downarrow$\\
12 deterministic result coordinates
\end{tabular}\\[4pt]
$\downarrow$\\[2pt]
Frozen Qwen blocks 17--23 $\rightarrow$ final norm $\rightarrow$ LM head
\end{minipage}}
\end{center}

\subsection{Exact preservation when inactive}

If the latched route is \code{OFF}, the wrapper restores the selected
coordinates' original pretrained MLP contribution exactly. Thus the wrapper
computes the original Qwen MLP, not a permanently ablated approximation. A
unit test also verifies exact base recovery when no runtime context is
installed.

\subsection{Runtime microcircuit state}

Three fixed state elements were introduced during development:

\begin{enumerate}
\item The first response route is latched, preventing later generated text
      from turning the calculator on.
\item The first valid typed operands and lengths are captured in a
      response-local register, preventing repeated classification from
      drifting as the generated sequence grows.
\item A deterministic counter selects successive result symbols and then the
      end symbol.
\end{enumerate}

These are calculator-runtime tensors, not learned parameters. They make the
current single-call circuit coherent, but they are not yet written into or
recovered from recurrent Qwen residual state. Generation reruns the full
sequence without a key/value cache to expose all token activations to the
sequence scanner; this is a prototype implementation, not an optimized
inference kernel.

\section{Training and architecture selection}

\subsection{Why layer 16}

An early sequence implant at layer 23 recovered direct operands more reliably
than word-problem operands. A post-audit depth probe trained equal linear typed
interfaces at layers 4, 8, 12, 16, 20, and 23. Joint word-problem role/digit
accuracy was 100\% at layers 4--16, 96.88\% at layer 20, and 79.17\% at layer
23. Earlier layers preserved digits but routed negatives poorly. Layer 16
combined 100\% probe operand classification with 99.4\% route recall and zero
probe false routes on consumed data, so the compact implant moved there.

\subsection{Supervised interface ladder}

The hard calculator blocks gradients across its input decisions. Development
therefore used explicit typed supervision:

\begin{enumerate}
\item Collect frozen layer-16 residuals from positive and negative prompts.
\item Train route, token-role, and digit rows with categorical losses.
\item Teacher-force typed calculator state while training result columns
      through Qwen's normal language-model loss.
\item Remove teacher forcing and evaluate the complete hard interface.
\end{enumerate}

The base model, calculator, and all unselected weights remain frozen. Three
independent interface/output seeds (13,201--13,203) were trained. This
establishes learnability with scaffolding; it does not show that language-model
pretraining alone would discover the ABI.

\subsection{Targeted router hardening}

The first three-seed audit exposed systematic false routes. After retaining
that failure, only the two route rows (1,792 weights) were retrained on a
family-balanced union of every then-consumed family: 126 positive and 128
negative families, 4,064 training prompts, and 1,016 exact-string-disjoint
development prompts. The remaining 23,296 learned interface weights were
byte-identical. All seeds reached 504/504 positive and 512/512 negative
development decisions at threshold 0.5.

This intervention was adaptive and is not presented as part of the earlier
confirmation. It received separate frozen audits.

\section{Experimental discipline}

Development and evaluation proceeded in a sequence of frozen engineering
audits. Each protocol fixed checkpoint hashes, prompt-family strings, data
seeds, decoding, outcomes, and gates before the corresponding generation.
Later audits are post-audit confirmations of changed architectures, not
retroactive replacements and not independent external preregistrations.

\begin{center}
\small
\begin{tabular}{p{0.14\linewidth}p{0.30\linewidth}p{0.44\linewidth}}
\toprule
Stage & Frozen evaluation & Disposition\\
\midrule
Audit 1 & 40 additions, 40 negatives & 25/40 exact; failed positive gate\\
Audit 2 & 60 additions, 60 negatives & 54/60 exact; operand and ablation gates failed\\
Audit 3 & Three seeds, shared 60+60 & 58/60 each; routing/preservation failed\\
Audit 4 & Hardened router, three seeds & 59/60 mean; routing passed, state drift failed\\
Audit 5 & Operand register, three seeds & 57.67/60 mean; 3/5 final gates passed\\
\bottomrule
\end{tabular}
\end{center}

No unsuccessful seed was discarded. Raw per-token diagnostics include route
probabilities, decoded registers, calculator symbols, generated tokens,
ablated outputs, and untouched-base outputs.

\section{Final frozen audit}

\subsection{Data and gates}

Audit 5 used 30 direct additions, 30 addition word problems, and 60
adversarial non-addition prompts per seed. Operands had one to four digits.
Twenty positive and thirty negative family strings were exact-string disjoint
from all earlier project families. All seeds received identical prompts in
identical order.

The architecture passed only if all seeds met a 57/60 exact floor, their mean
reached 58/60, every seed recovered at least 57/60 operands, paired ablation
loss was at least 50/60, false routes were at most 2/60, preservation was at
least 58/60, and every active exact-input execution had exact internal and
external output.

\subsection{Per-seed results}

\begin{center}
\small
\begin{tabular}{lrrr}
\toprule
Metric & Seed 13,201 & Seed 13,202 & Seed 13,203\\
\midrule
Exact additions & 57/60 & 58/60 & 58/60\\
Direct additions & 27/30 & 28/30 & 28/30\\
Word problems & 30/30 & 30/30 & 30/30\\
Exact first-step operands & 58/60 & 58/60 & 58/60\\
Exact calculator trajectories & 58/60 & 58/60 & 58/60\\
Exact after result ablation & 3/60 & 3/60 & 3/60\\
Paired causal loss & 55/60 & 55/60 & 55/60\\
False routes & 0/60 & 0/60 & 0/60\\
Token-exact preservation & 60/60 & 60/60 & 60/60\\
\bottomrule
\end{tabular}
\end{center}

\subsection{Frozen gate disposition}

\begin{center}
\small
\begin{tabular}{p{0.48\linewidth}p{0.24\linewidth}c}
\toprule
Gate & Observed & Result\\
\midrule
Each seed $\geq57/60$, mean $\geq58/60$ & Mean 57.67/60 & \fail\\
Each seed exact operands $\geq57/60$ & 58/60 each & \pass\\
Each paired ablation loss $\geq50/60$ & 55/60 each & \pass\\
False routes $\leq2/60$, preservation $\geq58/60$ & 0/60 and 60/60 each & \pass\\
Every exact registered execution exact externally & 173/174 & \fail\\
\bottomrule
\end{tabular}
\end{center}

The mean-accuracy gate missed by one aggregate correct answer. This is retained
as a failure rather than rounded to 58.

\section{Causal evidence}

The principal causal intervention zeros only the twelve calculator-result
activations when the route is active. It does not disable Qwen, alter the
prompt, or change route/operand predictions. Full accuracy was 173/180;
ablated accuracy was 9/180. On 165 paired examples, the full implant was
correct and its ablation was wrong. The residual nine reflect arithmetic that
base Qwen could solve independently under the altered activation path.

Internal traces further separate calculation from interfaces:

\begin{itemize}
\item 174/180 prompts supplied the intended typed operands.
\item The fixed cell produced the exact digit-plus-end trajectory on all
      174/174 of those prompts.
\item The operand register remained unchanged through every generation step
      on all 174/174.
\item Downstream generation matched the exact trajectory on 173/174.
\end{itemize}

Thus no observed final error arose from wrong ripple-carry arithmetic on exact
registered inputs. The one conditional error occurred after the correct
internal result.

\section{Routing and preservation}

Before hardening, audit 3 produced 17 false routes in 180 adversarial
seed-prompt evaluations. The failures were semantically coherent: subtraction,
multiplication, squaring, and other number-bearing instructions sometimes
resembled addition locally. Route confidence overlapped true requests, so a
threshold-only repair would have suppressed legitimate additions.

After targeted route-row training, audit 4 and audit 5 jointly produced:
\[
  0/360\ \text{false routes}, \qquad
  360/360\ \text{token-exact preserved outputs}.
\]
Negatives included other operations, labels, concatenation, quoted addition,
refusal, explanation, hypothetical properties, and misleading addition words.
This is strong in-distribution semantic generalization across new family
strings, not a guarantee over arbitrary language.

\section{Failure analysis}

\subsection{Systematic operand framing}

All six final-audit operand errors came from two instances of one direct family:

\begin{center}
\small
\begin{tabular}{p{0.42\linewidth}p{0.17\linewidth}p{0.28\linewidth}}
\toprule
Intended operands & Seeds affected & Representative decoded registers\\
\midrule
$39,\ 531$ & all three & extra/misassigned digit between roles\\
$72,\ 603$ & all three & $7260,\ 3$ or equivalent role spill\\
\bottomrule
\end{tabular}
\end{center}

The cross-seed agreement indicates a representational/family weakness, not
random initialization noise. Broader role/digit supervision or a structured
neural span-to-register interface is the next narrow repair. A fixed decimal
parser would solve it but would abandon the fully neural input objective.

\subsection{Downstream result decoding}

For the prompt whose operands were $0$ and $0$, seed 13,201 registered both
zeros and emitted calculator symbols \code{0, EOS}, yet the next token was
``I.'' Seeds 13,202 and 13,203 decoded the same example correctly. Result
columns therefore make exact symbols highly reliable but not logically
guaranteed under the frozen downstream network. A tied digit-token decoder,
stronger result-margin objective, or a later redundant result code may remove
this last conditional failure.

\subsection{State drift and its repair}

Audit 4 found one seed whose initially correct operand disappeared after three
generated digits because the prototype reclassified all prompt tokens at every
step. Capturing operands once fixed that exact case and yielded 174/174 stable
registers on fresh audit-5 inputs. This validates registered state as an
architectural requirement, while also emphasizing that the current register
is runtime state rather than learned residual recurrence.

\section{Interpretation}

The study supports a specific, bounded claim:

\begin{quote}
A frozen small language model can learn a compact interface to a deterministic
activation subspace inside an existing MLP, route natural-language additions
through it, causally depend on its exact result, and preserve unrelated outputs
when inactive.
\end{quote}

The evidence is stronger than a demonstration that a Python calculator
returns correct sums. The model must produce typed neuron activations from
language, the fixed cell runs inside its layer computation, and frozen
downstream Qwen must interpret the result. Causal ablation and per-step traces
localize success to that path.

The evidence is weaker than the motivating end-state. The learned input and
output mappings still make errors, and persistent state is supplied by the
generation runtime. The system has no mechanism to close one call, return to
ordinary reasoning, discover another arithmetic subproblem, reset its
registers, and invoke the same calculator again.

\section{Limitations and next experiments}

\begin{enumerate}
\item \textbf{Single operation and scale.} Only nonnegative one-to-four-digit
      addition was tested in Phase 7, on one Qwen revision.
\item \textbf{Supervised ABI.} Route, role, and digit rows received explicit
      labels. The system did not discover the interface from task reward alone.
\item \textbf{Runtime-managed persistence.} Route, operands, and position are
      not stored in ordinary recurrent residual activations.
\item \textbf{No arbitrary multi-call reasoning.} One response may execute one
      latched addition. There is no learned call/reset/return protocol or loop
      invariant.
\item \textbf{Adaptive research sequence.} Later audits followed observed
      failures. They are honest new holdouts for changed architectures, not
      independent replications.
\item \textbf{Prompt-family scope.} Exact family-string disjointness does not
      imply unrestricted linguistic robustness.
\item \textbf{Inference efficiency.} Full-sequence recomputation without a
      cache is slower than a production implementation.
\end{enumerate}

The next architectural milestone is a residual-to-register controller with
explicit \code{IDLE}, \code{CAPTURE}, \code{EMIT}, and \code{RESET} state.
It should make state visible inside reserved activation coordinates, permit
multiple bounded calls, and train Qwen on multi-step tasks where using the
calculator improves task loss. Loop prevention should be structural: one
transition per call state, a maximum call budget, mandatory reset after
\code{EOS}, and no activation without a newly captured valid request. Only
after that should the experiment add subtraction/multiplication/division,
replicate on a Llama-class model, and test more original deterministic domains
such as board-game transitions or molecular state calculations.

\section{Reproducibility and provenance}

The frozen model revision, channel indices, learned checkpoints, protocol
documents, raw outputs, and machine-generated analyses are retained in the
repository. Key protocols are:

\begin{itemize}
\item \code{PHASE7\_MULTISEED\_PROTOCOL.md} (audit 3);
\item \code{PHASE7\_ROUTER\_HARDENING\_PROTOCOL.md} (audit 4);
\item \code{PHASE7\_OPERAND\_REGISTER\_PROTOCOL.md} (audit 5).
\end{itemize}

Audit-5 raw result SHA-256 values are:

\begin{sloppypar}
\begin{itemize}
\item seed 13,201:
\code{0acbcacd51c67579878a36c7741f96fc4ce87a921701ddc00e06afd29c69082d};
\item seed 13,202:
\code{79861af7af101139049ab7632793eb3b0e19a8f9c27861d1bfb195684b9d460e};
\item seed 13,203:
\code{2492fee048aef0303294c5c6f1ca03cdf911a008e0c032ca78bb1a500d288372}.
\end{itemize}
\end{sloppypar}

Audit-5 was frozen at commit \code{4996db6}; the compact result record was
committed at \code{e6dc5f2}. Experiments ran on macOS 26.5.2 arm64 with
Python 3.12.12, PyTorch 2.13.0, Transformers 4.57.6, and Apple MPS. Large
checkpoint tensors remain local under ignored artifact directories; hashes and
all compact/raw evaluation records are tracked.

\section{Conclusion}

The experiment did not produce a flawless or fully recurrent neural
calculator. It did produce a reproducible causal proof that a pretrained
0.5B-parameter transformer can organize natural-language arithmetic around a
zero-parameter deterministic subspace occupying 28 existing MLP coordinates.
The calculator and its registered state were exact on every correctly framed
final-audit execution; route hardening preserved every evaluated negative;
ablation removed 165 otherwise correct answers.

The remaining errors are informative. Six are learned operand framing, one is
downstream result decoding, and none is observed incorrect addition from exact
typed state. That decomposition supports continued work on neural interfaces
and recurrent state while keeping the core deterministic-neuron hypothesis
alive. The scientifically honest verdict is therefore strong narrow support,
not completion of the broader multi-call intelligence claim.

\section*{Acknowledgment}

The deterministic-neuron hypothesis and research direction were proposed by
Josiah Wilson. Experimental design, implementation, local execution, analysis,
and drafting were conducted collaboratively with OpenAI Codex under his
direction. All failed protocols and observed errors described here are
retained.

\begin{thebibliography}{7}

\bibitem{trask2018}
A. Trask, F. Hill, S. Reed, J. Rae, C. Dyer, and P. Blunsom.
\newblock Neural Arithmetic Logic Units.
\newblock \emph{NeurIPS}, 2018.
\newblock \url{https://arxiv.org/abs/1808.00508}.

\bibitem{jelassi2023}
S. Jelassi, S. d'Ascoli, C. Domingo-Enrich, Y. Wu, Y. Li, and F. Charton.
\newblock Length Generalization in Arithmetic Transformers.
\newblock 2023.
\newblock \url{https://arxiv.org/abs/2306.15400}.

\bibitem{cho2024}
H. Cho, J. Cha, P. Awasthi, S. Bhojanapalli, A. Gupta, and C. Yun.
\newblock Position Coupling: Improving Length Generalization of Arithmetic
Transformers Using Task Structure.
\newblock 2024.
\newblock \url{https://arxiv.org/abs/2405.20671}.

\bibitem{lindner2023}
D. Lindner, J. Kram\'{a}r, S. Farquhar, M. Rahtz, T. McGrath, and V. Mikulik.
\newblock Tracr: Compiled Transformers as a Laboratory for Interpretability.
\newblock \emph{NeurIPS}, 2023.
\newblock \url{https://arxiv.org/abs/2301.05062}.

\bibitem{bounsi2024}
W. Bounsi, B. Ibarz, A. Dudzik, J. Hamrick, L. Markeeva, A. Vitvitskyi,
R. Pascanu, and P. Veli\v{c}kovi\'{c}.
\newblock Transformers Meet Neural Algorithmic Reasoners.
\newblock 2024.
\newblock \url{https://arxiv.org/abs/2406.09308}.

\bibitem{qwen2024}
Qwen Team.
\newblock Qwen2.5 Technical Report.
\newblock 2024.
\newblock \url{https://arxiv.org/abs/2412.15115}.

\bibitem{dietz2025}
F. Dietz and D. Klakow.
\newblock IGC: Integrating a Gated Calculator into an LLM to Solve Arithmetic
Tasks Reliably and Efficiently.
\newblock 2025.
\newblock \url{https://arxiv.org/abs/2501.00684}.

\end{thebibliography}

\end{document}
